\begin{table}[!t]
    \caption{
        评估中使用的模型规格。
    }
    \label{table:model-params}

    \footnotesize
    \begin{tabular}{lccccc}
    \toprule
    \textbf{名称} & \textbf{\makecell{层数}} & \textbf{\makecell{嵌入\\维度}} & \textbf{\makecell{FFN 隐藏\\维度}} & \textbf{\makecell{注意力\\头数}} & \textbf{\makecell{注意力\\组数}} \\
    \midrule
    ViT 5B~\cite{vit22b-arxiv23} & 63 & 1792 & 15360 & 16 & 16 \\
    ViT 22B~\cite{vit22b-arxiv23} & 48 & 6144 & 24576 & 48 & 48 \\
    \midrule
    Llama3 8B~\cite{llama3-arxiv24} & 32 & 4096 & 14336 & 32 & 8 \\
    Qwen2 32B~\cite{qwen2-5-arxiv25} & 64 & 5120 & 27648 & 40 & 8 \\
    Qwen2 72B~\cite{qwen2-5-arxiv25} & 80 & 8192 & 29568 & 64 & 8 \\
    \midrule
    DiT 5B~\cite{moviegen-arxiv24} & 28 & 3584 & 10240 & 28 & 28 \\
    DiT 30B~\cite{moviegen-arxiv24} & 48 & 6144 & 24576 & 48 & 48 \\
    \bottomrule
    \end{tabular}
\end{table}

\begin{table}[!t]
    \caption{
        评估中使用的模型组合。
    }
    \label{table:model-setup}

    \small
    \begin{tabular}{llccc}
    \toprule
    \textbf{名称} & \textbf{模型配置} & \textbf{TP} & \textbf{PP} & \textbf{GPU 数} \\
    \midrule
    VLM-S & ViT 5B + Llama3 8B & 4 & 4 & 16 \\
    VLM-M & ViT 5B + Qwen2 32B & 8 & 4 & 32 \\
    VLM-L & ViT 22B + Qwen2 72B & 8 & 8 & 64 \\
    \midrule
    T2V-S & Llama3 8B + DiT 5B & 4 & 4 & 16 \\
    T2V-L & Qwen2 32B + DiT 30B & 8 & 8 & 64 \\
    \bottomrule
    \end{tabular}
\end{table}

\section{实现}
\label{section:implementation}

\sys{} 以 Megatron-LM~\cite{megatron-lm-arxiv20} 的扩展形式实现，
其中包含一个与 Megatron-LM 运行时协同工作的中央规划器。
规划器反复预取训练元数据、搜索流水线调度，
并将优化后的调度部署到 GPU 集群。

\subsection{训练仿真器}
\label{section:multimodal-training-simulator}

\sys{} 使用算子级分析建模进行性能预测。
仿真器构建包含两类节点的有向无环图（DAG）：
(1) 表示底层 GPU 操作（例如矩阵乘法和集合通信）的算子节点；
(2) 对应数据缓冲区（例如模型参数）的张量节点。
每个节点被分配到特定设备（GPU、CPU 或 NVLink），并通过依赖边相连。

在延迟估计中，每个算子节点由以下量刻画：
浮点运算次数 $N_\mathrm{fop}$、
以字节计的内存访问量 $N_\mathrm{mem}$，
以及以字节计的网络传输量 $N_\mathrm{net}$。
给定设备能力（以 FLOPS 计的计算能力 $F$、
以字节/秒计的内存带宽 $B_\mathrm{mem}$，
以及以字节/秒计的网络吞吐量 $B_\mathrm{net}$），延迟按下式计算：
$
\max\{
    \alpha_\mathrm{fop}N_\mathrm{fop}/F,
    \alpha_\mathrm{mem}N_\mathrm{mem}/B_\mathrm{mem},
    \alpha_\mathrm{net}N_\mathrm{net}/B_\mathrm{net}
\}
$，
其中 $\alpha_\mathrm{fop}$、$\alpha_\mathrm{mem}$ 和 $\alpha_\mathrm{net}$ 是效率缩放因子。
用户也可以使用自定义估计规则覆盖该成本模型。

仿真器按拓扑顺序填充算子时间戳，
随后通过检查相关算子节点确定张量生命周期。
利用这些生命周期，可以构建内存消耗时间线，
从而估计各设备的峰值内存使用量。

\subsection{并行调度搜索}

\sys{} 在 CPU 核心上执行训练仿真器和搜索算法。
它首先使用预取的元数据，对各个微批次计算图进行仿真，以获取流水线阶段延迟。
调度搜索期间，多个工作进程并行探索调度空间，并共享全局 MCTS 搜索统计信息。
每轮搜索结束后，这些工作进程通过互斥锁保护的操作，以原子方式更新全局 MCTS 统计信息。
由于工作进程每次同步之间会执行多次回放，从而摊薄同步开销，因此竞争始终很小。
为避免干扰正常训练进程，\sys{} 将搜索工作进程数量限制为可用 CPU 核心的至多 50\%
（例如一台配有 8 块 H800 GPU 的机器上使用 64 个核心）。

\subsection{执行计划部署}

将仿真的流水线调度部署到 GPU 集群，需要把它们编译为指定计算和通信模式的物理执行计划。

\heading{调度编译}
\sys{} 沿用 DynaPipe~\cite{dynapipe-eurosys24}，定义构成流水线执行计划的动作序列：
\begin{itemize}
    \item 流水线阶段被转换为 \texttt{fw\_stage}/\texttt{bw\_stage} 动作，并采用\S\ref{section:memory-optimization}给出的优化策略。
    \item 点到点（P2P）通信使用与计算重叠的异步内核（\texttt{isend}、\texttt{irecv}）。
    \item 根据仿真时间线插入同步动作（\texttt{wait\_isend}、\newline\texttt{wait\_irecv}）。
\end{itemize}

\heading{运行时修改}
我们扩展了 Megatron-LM 的 \texttt{schedules} 模块以支持动态计划。
每个流水线工作进程通过 RPC 从中央规划器接收一个动作列表，
并按顺序执行其中的动作。
连续的 P2P 内核会被组合成批处理操作，以提高效率。
